% !TeX root = ../../../main.tex
\subsection{Discussion and Conclusion}

Our findings highlight the superiority of encoder-based models over decoder-based LLMs for event forecasting, emphasizing the limitations of autoregressive models in structured prediction tasks.
Among encoder models, \textit{BioClinical-ModernBERT-base} consistently performs best, achieving superior concordance (0.677 vs 0.646) and better external generalization despite slightly lower internal F1 scores than general-purpose \textit{ModernBERT-base}.
MLP-head fine-tuning excels in F1 scores while fine-tuned masking models achieve higher concordance, particularly in long-term forecasting.
This suggests that classification and ranking tasks benefit from distinct optimization strategies.

The consistent advantage of biomedically-pretrained models across tasks deserves special attention.
While general-purpose models achieve strong internal performance, biomedical pretraining provides crucial benefits for real-world deployment: \textit{BioClinical-ModernBERT} variants show the smallest performance degradation on external datasets and maintain higher zero-shot capabilities (0.246 F1 at 1h vs near-zero for standard \textit{BERT}/\textit{RoBERTa}).
This suggests that domain-specific pretraining enhances not just accuracy but also robustness and generalization—critical factors when considering clinical deployment where distribution shifts are common.

The gradual improvement of F1 scores over longer time horizons indicates that event patterns become more predictable over time, whereas short-term forecasting remains challenging due to higher variability.
Performance drops on external validation datasets highlight generalization challenges, though \textit{ModernBERT} models exhibit relative robustness.
For survival analysis, instruction-tuned LLMs like \textit{Llama-3.3-70B-Instruct} outperform traditional transformer baselines, with several models achieving peak concordance at early timepoints rather than with extended observation.

Our sensitivity analysis reveals several trade-offs: time-ordered training generally improves concordance while text-ordering can yield better F1 scores on externally-annotated datasets.
Robustness experiments show F1 scores degrade significantly beyond 60\% timestep dropout, while concordance remains stable, indicating event ranking is less sensitive to partial history than event classification. 

Additional insights into the forecasting results are presented in Appendix~\ref{ch07:apd:additional_insights}.
Specifically, we summarize two observations that clarify model behavior and apparent anomalies: (i) why decoders fare better on survival but not short-horizon forecasting, and (ii) interpreting very low $F_{1}$ scores and external validation drops.

\textbf{Methodological Contributions and Potential Impact.} Our framework demonstrates how narrative clinical texts can be systematically converted into structured temporal representations for forecasting tasks. The ability to extract temporally structured insights from unstructured clinical text could potentially support clinical decision-making, particularly in settings with limited structured data or specialized expertise. 

Importantly, while consumers frequently consult LLMs about health risk via prompting in chats, this study demonstrates that the prompting approach performs substantially worse in risk prediction than alternative approaches, at least with respect to precision/recall/F1, with prompted LLMs achieving at best 0.460 F1 at 24h compared to 0.653 for fine-tuned encoders.
Our study highlights several alternative approaches and characterizes their relative performance strengths and weaknesses.
In addition, our finding that concordance remains stable with 60\% missing context suggests robustness to incomplete documentation scenarios, which highlights the degradation pattern in performance as it relates to context availability. 

Our approach also demonstrates that our language model systems can reliably capture the temporal reasoning that clinicians use in practice based on their clinical documentation.
However, additional work would be required to validate performance on real-time clinical data and demonstrate measurable impact on patient outcomes in deployment.
Further discussion on societal impact of our work is in Appendix \ref{ch07:apd:broader_impact}.

\textbf{Limitations and Future Directions.} 
Our study has important limitations to consider.
First, our pipeline relies on case reports from the PubMed Open Access (PMOA) corpus.
Because these reports often highlight rare or atypical presentations and differ from routine clinical notes (e.g., progress notes, discharge summaries), generalizability to real-world health-system text may be limited.
However, case reports provide rich, temporally structured narratives with explicit clinical rationale, making them well suited for evaluating model performance under sparse documentation.
Accordingly, we present a general framework for textual time series forecasting that uses case reports as an interpretable, temporally rich foundation for method development, with extensions to real-world corpora such as MIMIC discharge summaries.
See Appendix~\ref{ch07:apd:beyond_pmoa} for a detailed discussion on the generalizability of our pipeline beyond case report details, as well as comparisons with other clinical sources.

Second, while we focused on sepsis due to its clinical relevance and prevalence in case reports, our framework is fundamentally disease-agnostic.
Preliminary results from over 125K PMOA case reports across various diagnostic conditions show promising generalization, with future work planned on broader diagnostic categories and real-world corpora such as MIMIC-IV discharge summaries.

Finally, despite mitigation steps (timestamp extraction, temporal masking, external test sets), pretrained language models may still have been exposed to PMOA content.
Although they were not trained on our derived \emph{(event, time)} sequences or temporally reordered data, the original narratives could appear in pretraining.
Our pipeline applies temporal reordering and masking of future events beyond the prediction cutoff, enforcing a causal structure absent from the originals and yielding a distinct setup that requires reasoning over partial, temporally grounded inputs. 
% To assess any pretraining exposure, we include \textit{RedPajama-INCITE-7B-Instruct}, an open-source model with a documented PubMed-free corpus; its performance (0.352 F1 at 24h; concordance 0.618), comparable to decoders potentially exposed to PubMed, supports that the T2S2 task is sufficiently distinct from raw PMOA text and mitigates causal leakage concerns.
While we included \textit{RedPajama-INCITE-7B-Instruct}, an open-source model with a documented PubMed-free corpus for pretraining, and saw its performance (0.352 F1 at 24h; concordance 0.618) was comparable to decoders potentially exposed to PubMed, other strategies supporting the argument that the \textit{T2S2} task is sufficiently distinct from raw PMOA text and mitigates causal leakage concerns could be examined.
See Appendix~\ref{ch07:apd:open_source} for a detailed discussion of how open-source models mitigate potential data leakage, as well as additional forecasting results on 115 case reports published after the pretraining cutoffs of encoder-based models (post-2020).

%Finally, our analysis is limited by the modest size of the high-cost manual annotations collected as a gold standard to calibrate and validate our LLM extraction pipeline. Expanding gold-standard validation across diverse diagnoses could assess the generalizability of our approach and aid in evaluating model robustness, and is a direction for future work.

%%%%%%%%%%%%%%%%%%%%%%%%%%%%%%%%%%%%%%%%%%%%%%%%%%%%
%%%%%%%%%%%%%%%%%%%%%%%%%%%%%%%%%%%%%%%%%%%%%%%%%%%

